\section{Additional method details}
\label{app:method}

This appendix records implementation details inherited from text-only CoBit
\citep{method:cobit}. The main paper focuses on the multimodal representation and partial-denoising
objective used in this work.

\subsection{Continuous bitstream diffusion}

CoBit uses a variance-exploding Gaussian diffusion with, by default,
$\sigma\in[0.002,80]$ (CoBit-MLS trains with $\sigma_\text{min}=0.01$), data centre $1/2$ and $\sigma_{\mathrm{data}}=1/2$. The analytic posterior
logit of an isolated bit under a uniform prior,
$(\vx_\sigma-\tfrac12)/\sigma^2$, is clipped to $\pm30$ and added to a learned contextual
residual. The network therefore concentrates on dependencies between bits rather than relearning
the local binary denoising problem.

The entropy-derived schedule estimates the conditional entropy rate from the unweighted
per-example error, $\hat h(\sigma)\propto e(\sigma)/\sigma^2$. Estimates are stored in a FIFO buffer
of $8\times10^5$ examples and 128 log-spaced bins. The training proposal and sampling grid use
$q(\sigma)\propto g(\sigma;c,n)\hat h(\sigma)^{1/2}$, where
$g(\sigma;c,n)=\sigma^n/(\sigma^n+c^n)$ with $c=0.1$ and $n=3$. Training begins with the EDM
log-normal proposal ($P_{\mathrm{mean}}=-1.2$, $P_{\mathrm{std}}=1.2$), interpolates to $q$ over
steps $4$--$5\times10^4$ by default (CoBit-LibriTTS warms up over $1$--$1.3\times10^4$ steps), and refreshes the estimate every $2\times10^3$ steps.

For self-conditioning, the denoiser receives a second input $\hat\vx_0$. With probability
$p_{\mathrm{sc}}=0.5$, a detached no-gradient pass supplies this estimate during training;
otherwise it is zero. At inference, \texttt{carry} mode passes the previous prediction to the next
step without another evaluation; the \texttt{refined} mode, used for the protein evaluations, recomputes the estimate
at the new state and doubles the cost per step.

Sampling integrates the probability-flow ODE with a first-order DDIM/Euler step on the
entropy-derived grid, optionally with EDM-style churn \citep{method:ddim,image:edm}. Final
probabilities are thresholded at $1/2$ and regrouped into codes. The Sequence Diffusion Transformer
patches the $m$ bits at each position into one trunk token, uses RoPE, SwiGLU and adaLN-zero noise
conditioning, and combines the trunk output with the raw noisy bits in a light skip-connection
head.

\subsection{Representation and partial denoising}

Each frozen tokenizer maps its modality to integer codes and decodes predicted codes after
sampling. Per-modality vocabularies occupy disjoint ranges of one integer code space together with
padding and structural markers. The bit width $m$ covers their union, so namespace information is
part of the same binary word as token content. Markers remain clean throughout training and
sampling. Fixed token budgets are truncated or padded, and two-dimensional grids are flattened in
the order specified by the experiment.

Training draws a conditioning regime for each example. The mask $M$ marks clean conditioning bits
and the loss is evaluated only on $1-M$. When self-conditioning is active, clean conditioning bits
are also copied into $\hat\vx_0$.

% For classifier-free guidance, conditioning segments are replaced
% by the null value $1/2$ with probability $p_{\mathrm{uncond}}=0.1$ while markers remain protected.
% Examples whose conditioning was dropped are excluded from the entropy-rate buffer so that the
% estimated schedule reflects the conditional tasks.

At inference, the task fixes $M$. Conditioned positions are clamped before every denoiser call and
their ODE velocity is set to zero. Guidance combines conditional and unconditional predictions in
probability space as $D_g=D_u+w(D_c-D_u)$. The two branches are evaluated as one batch and carry
separate self-conditioning estimates. Joint generation uses one evaluation per step; guided
conditional generation uses two.
